\documentclass[../../../include/open-logic-section]{subfiles}

\begin{document}

\olfileid{sth}{ordinals}{vn}	
\olsection[Von Neumann's Construction]{von Neumann（冯·诺伊曼）的序数构造}

\olref[sth][ordinals][iso]{thm:woalwayscomparable} 引发了一个想法。
我们可以引入一些称为\emph{序型}的对象，来充当良序的代表。
用 $\ordtype{A, <}$ 表示良序 $\tuple{A, <}$ 的序型，我们希望确立以下两条原则：
\begin{align*}
	\ordtype{A, <} = \ordtype{B, \lessdot} & 
	\text{ 当且仅当 } \ordeq{\tuple{A, <}}{\tuple{B, \lessdot}}\\
	\ordtype{A, <} < \ordtype{B, \lessdot}&
	\text{ 当且仅当 }\ordeq{\tuple{A, <}}{\tuple{B_b, \lessdot_b}}\text{ 对于某个 }b \in B
\end{align*}
此外，我们可能还希望将序型\emph{作为某类集合}引入，就像我们可以将自然数作为某类集合引入一样。

实现这一目标最常用的方法——也是我们将遵循的路径——是通过某些\emph{典范的}良序集来定义这些序型。
这些典范集合最早由 von Neumann 引入：

\begin{defn}
集合 $A$ 是\emph{传递的}，当且仅当 $(\forall x \in A)x \subseteq A$。
$A$ 是\emph{序数}，当且仅当 $A$ 是传递的且被 $\in$ 良序化。
\end{defn}

\noindent
在下文中，我们将用希腊字母表示序数。
由定义立得，若 $\alpha$ 为序数，则 $\tuple{\alpha, \in_\alpha}$ 为良序，其中 $\in_\alpha =
\Setabs{\tuple{x, y} \in \alpha^2}{x \in y}$。
所以，稍加滥用记号，我们可以直接说 $\alpha$ \emph{本身}就是一个良序。

以下是我们的前几个序数：
\[
	\emptyset, \{\emptyset\}, 
	\{\emptyset, \{\emptyset\}\}, 
	\{\emptyset, \{\emptyset\}, \{\emptyset, \{\emptyset\}\}\}, \ldots
\]
你会注意到，这些正是我们在无穷公理中，即在 von Neumann 对 $\omega$ 的定义中（参见 \olref[sth][z][infinity-again]{sec}）最初遇到的那几个序数。
这并非巧合。von Neumann 的序数定义将自然数视为序数，但也允许存在超限序数。

和往常一样，我们现在可以问：这些\emph{就是}序数吗？还是说 von Neumann 仅仅给了我们一些可以\emph{当作}序数来处理的集合？
关于这个问题的讨论，与我们之前在 \olref[sfr][rel][ref]{sec}、
\olref[sfr][arith][ref]{sec}、
\olref[sfr][infinite][dedekindsproof]{sec} 以及
\olref[sth][z][nat]{sec} 中进行的讨论类似，因此我们就不在此赘述了。
在后续行文中，我们将简单地用“序数”来指代“冯·诺伊曼序数”。

\end{document}